\documentclass[11pt]{article}
\usepackage[letterpaper,margin=0.92in]{geometry}
\usepackage{amsmath,amssymb,amsthm}
\usepackage[hidelinks]{hyperref}
\newcommand{\Rea}{\operatorname{Re}}
\newcommand{\supp}{\operatorname{supp}}
\title{Close zeros, spectral lifts, and a test for the cubic clock}
\author{}\date{}
\begin{document}
\maketitle
\vspace{-2.8em}

\section*{1. A uniform bound in the spectral trace proof}
In the supplied version of Alessandro Maria Masullo's \emph{The Music of the
Primes}, the proof of Theorem 10.4.2.8 requires lifts of the biorthogonal
vectors with norms bounded by a single $C+\varepsilon$ in (10.4.2.24),
p.~463. Proposition 10.2.2.2, cited there, concerns the norms of the
evaluation functionals $\eta$ on pp.~446--447. The following specializes
the issue to the Riemann zeta function.

Take the trivial character, $a=1/2$, a fixed strip $I=[-\alpha,\alpha]$
of positive width and a fixed Sobolev index $s>1$. In the centered
coordinates of (10.2.2.4), the evaluation at a critical-line zero
$1/2+it$ is $\eta_t(u)=e^{itu}$, with $u=\log|g|$. The compact norm-one
factor contributes a constant Haar volume. Equation (10.2.2.8) and its
proof give the dual norm on this component with density
\[
 w(u)=e^{-2\alpha |u|}(1+u^2)^{-s/2}.
\]
In particular $\|\eta_t\|$ is constant. Yet for any two ordinates $t,t'$,
\[
 \|\eta_t-\eta_{t'}\|\le C_1|t-t'|,\qquad
 C_1^2=\int_{\mathbb R}u^2w(u)\,du
\]
up to the same Haar normalization. If the uncentered convention of
(10.4.2.22) is used, the extra fixed factor $e^{au}$ merely changes the
constant on choosing $\alpha>|a|$, an allowed strip.
For any lift $\psi_t$ biorthogonal to the evaluations, the pairings are
one at $t$ and zero at $t'\ne t$. Consequently
\[
 1=|\langle\psi_t,\eta_t-\eta_{t'}\rangle|
 \le C_1\|\psi_t\|\,|t-t'|.
\]
Conrey's unconditional positive proportion of simple critical-line zeros
and the Riemann--von Mangoldt count give $\gg T\log T$ distinct such
ordinates through $T$. Hence some gaps are $O(1/\log T)$.
The uniform bound on the lifts is impossible. The estimate applies to
every lift, since the evaluations factor through the quotient. It leaves
open bounds for sums of projections or for grouped zeros.

\paragraph{An amplification.}
In centered coordinates fix $0<r<\alpha$ and an interval length
$\ell=r/2$. Pigeonholing up to $T$ gives an interval of length $\ell$
containing $k+1\gg_\ell\log T$ distinct simple ordinates $t_0,\ldots,t_k$.
For $f(v)=\Rea\langle\psi_{t_0},\eta_v\rangle$, interpolation at
$t_1,\ldots,t_k$ gives $1\le\sup|f^{(k)}|\ell^k/k!$.
Since $|u|^k/k!\le r^{-k}e^{r|u|}$, the derivative is bounded by
$C_r\|\psi_{t_0}\|k!r^{-k}$. Thus
$\|\psi_{t_0}\|\ge C_r^{-1}2^k$: the largest required norm through
$T$ grows at least as a positive power of $T$. Orthogonal projection
onto a joint subspace avoids the requirement of separating each pair.
This calculation audits the particular estimate, not every trace identity.

\newpage
\section*{2. A single smooth test in a moving prime window}
The cubic $F=\mathbb Q[x]/(x^3-4x-1)$ has
$A(s)=\zeta_F(s)/\zeta(s)=L_k(s,\xi)$ for an order-three ideal-class
character of $k=\mathbb Q(\sqrt{229})$. Its completed $L$-function is
entire and self-dual. On this character component Masullo's
Theorem 12.5.1.1 (p.~519) asks the log-norm test to vanish at $u=0$.
The different condition in (13.3.1.2), p.~522, asks for two vanishing
Mellin moments because $[k:\mathbb Q]=2$. On this component they are
the ordinary moments of orders zero and one in $u$.

There is an explicit smooth test meeting both conditions and detecting
every off-line zero. Let $b$ be the density of the convolution of
uniform probabilities on $[-a_j,a_j]$, $a_j=2^{-j-2}$, $j\ge1$.
Its bilateral Laplace transform is
\[
 B(z)=\prod_{j\ge1}\frac{\sinh(a_jz)}{a_jz}.
\]
The product converges locally uniformly, and is nonzero for $\Rea z\ne0$
because $\sum_j a_j^2<\infty$ and each factor is nonzero there.
On any bounded vertical strip the first $m$ factors give decay of order
$m$, and the remaining factors have product bounded by
$\exp(|\Rea z|\sum_j a_j)$. Fourier inversion therefore gives a smooth
density supported on $[-1/4,1/4]$.

Set $g=b'$, $\widetilde g(u)=g(-u)$, $r=g*\widetilde g$. Its transform
is $R(z)=-z^2 B(z)^2$. In particular $R(0)=R'(0)=0$, $R(it)\ge0$,
and $R$ has no off-line zeros. For $T>1/2$ form the test
\[
 r_T(u)=\tfrac12\big(r(u-T)+r(u+T)\big).
\]
It vanishes near the origin, and its transform $R(z)\cosh(Tz)$ still
has a double zero at zero. Multiplying by the conjugate class character
on the norm-one group isolates the desired component. In the ordinary
explicit formula the finite-prime terms involve only prime powers in
$[e^{T-1/2},e^{T+1/2}]$.

\paragraph{Bounded response criterion.}
Let $E(T)=W_A(r_T)$ in the centered Weil formula. Then $E(T)$ is bounded
for $T>1/2$ if and only if all nontrivial zeros of $A$ lie on the
critical line. Indeed, putting $z=\rho-1/2$ and counting multiplicities,
the ordinary explicit formula and symmetry give
\[
 E(T)=\sum_z R(z)e^{Tz},\qquad
 \int_0^\infty E(T)e^{-sT}\,dT=\sum_z\frac{R(z)}{s-z}
 \quad(\Rea s>1/2).
\]
Rapid decay permits termwise integration and local meromorphic continuation.
On the critical line the first sum is uniformly bounded. Conversely,
boundedness would make the integral holomorphic on $\Rea s>0$, while
at any off-line zero there the sum has a pole of residue
$\operatorname{ord}(\rho)\,R(z)\ne0$. The functional equation supplies
the other half of the strip. A uniform estimate for these moving tests
remains to be proved. The point here is that the two printed test
restrictions can both be met. Theorem 13.3.1.2 in the supplied PDF is
an unfilled statement, and (12.5.1.3) has a strip-dependent zero sum.

\newpage
\section*{3. Two checks on the accompanying reports}
\paragraph{A discrepancy in the oper counts.}
For the separate cubic $x^3-2$, the Gemini Round 3 table uses
$f=t^3-2$ and, at $\mu=1/2$, the equation
$f y''+(3/2)t^2 y'+B y=0$. One can test dormancy over
$\mathbb F_p(t)$ directly. Put
\[
 M=\begin{pmatrix}0&f\\-B&-3t^2/2\end{pmatrix},\qquad
 U_0=I,\quad U_{j+1}=fU_j'-jf'U_j+U_jM .
\]
Iteration gives $(y,y')^{(j)}=f^{-j}U_j(y,y')$ in column notation.
The restricted $p$th power of $d/dt$ vanishes, so a zero matrix $U_p$
is exactly the zero $p$-curvature test. Polynomial arithmetic gives
\[
\begin{array}{c|c|c}
 p & B\text{ for which }U_p=0 &\text{Gemini report for the same curve}\\ \hline
 7 & 3,5,6&\text{no values}\\
 23 &0&0,3,6\\
 37 &\text{no values}&\text{three values}
\end{array}
\]
These direct results agree with the equation in the Astra Round 3 report
after the change of accessory parameter $c=-4B$. For example at $p=7$,
$[t^6](t^3-2)^3=1\bmod7$ and $(-4B)^3=1$ gives exactly $3,5,6$.
The comparison for this equation makes no claim about the other radii.

\paragraph{The paraboloid set can recover the stronger density.}
The Fable tower report constructs a bounded-relation-free set with
$|T|=A^d$, $\max T<n/(4L)$ and
\[
 |T|/n\ge 2^{-d}\Gamma^{-1}n^{-2/(d+2)},\qquad
 \Gamma=4dL(2L)^d,\qquad A=\lfloor(n/\Gamma)^{1/(d+2)}\rfloor .
\]
It takes $d$ fixed. In the schedule $2\le L\ll_\varepsilon\log n$,
choose instead $d=\lfloor\sqrt{\log n/\log(4L)}\rfloor$.
Then $d\to\infty$, $\log\Gamma=o(\log n)$, and $A\ge2$ eventually.
Taking logarithms gives
\[
 \log(n/|T|)\le d\log2+\log\Gamma+\frac{2\log n}{d+2}
 \ll\sqrt{\log n\log L},
\]
with an absolute implied constant for large $n$. This supplies even the
exponential-density conjunct in the existing Lean statement
\texttt{exists\_thick\_set\_sizes} (whose lower bound also divides by $L$).
The independent rank condition $m^4\le f$ remains in
\texttt{StandingNumericalData.lean}, line 118. This is an integer-set
construction and a source-level comparison, without a new Lean build.

\paragraph{Sources.}
A. M. Masullo, \href{https://math216.wordpress.com/wp-content/uploads/2026/11/the-music-of-the-primes.pdf}
{\emph{The Music of the Primes}}, supplied 940-page version,
especially pp.~446--447, 463, 519 and 522.
J. B. Conrey, \href{https://doi.org/10.1515/crll.1989.399.1}
{\emph{More than two fifths of the zeros of the Riemann zeta function are
on the critical line}}, J. reine angew. Math. 399 (1989), 1--26.
The local comparisons refer to the two \texttt{wall1\_r3} and three
\texttt{kill\_tower} reports in the workspace.
\end{document}
